% PSALM manuscript: IEEE single-column journal style (arXiv preprint).
% Body is written in continuous journal prose: no bulleted or numbered lists.
% Sections are filled incrementally FROM findings (closure-contract ARTIFACTS
% layer), never from the plan. Every citation must resolve to a real work
% (see .cursor/rules/citation-integrity.mdc); arXiv:2605.12548 is fabricated and
% must never appear here.
\documentclass[journal,onecolumn,11pt]{IEEEtran}

\usepackage{graphicx}
\usepackage{amsmath,amssymb}
\usepackage{booktabs}
\usepackage{url}
\usepackage[hidelinks]{hyperref}

\begin{document}

\title{PSALM: P\=a\d{n}inian Structured Pretraining for Small Language Models}

\author{PSALM Project\thanks{Preprint. Code, data, and models:
\url{https://github.com/SharathSPhD/PSALM}; Hugging Face: \texttt{qbz506/psalm-*}.}}

\maketitle

\begin{abstract}
Small language models are conventionally made competent through internet-scale
text. This work asks whether linguistic structure can substitute for data scale.
We treat P\=a\d{n}ini's A\d{s}\d{t}\=adhy\=ay\={\i}, the most formally complete
generative grammar written for any human language, as an unbounded data engine
that emits well-formed Sanskrit annotated with gold k\=araka roles, and we use
this stream as a pre-pretraining signal for a small multilingual decoder. We test
three claims: that a P\=a\d{n}inian structural prior improves sample efficiency
and compositional generalization relative to a matched artificial-grammar control
(H1); that a Navya-Ny\=aya reasoning scaffold reduces fallacious inference and is
more sample-efficient to install on a grammar-structured base (H2); and that an
epistemic constraint kernel enforces validity by construction at inference (H3).
All experiments run on a single accelerator at fixed token budgets. We report
results with explicit uncertainty and matched ablations, and we are deliberate
about what the approach does not claim. % Results sentence is filled at closure.
\end{abstract}

\section{Introduction}
% Filled from findings. Motivate structure-vs-scale; state the thesis and the
% three hypotheses in prose. No lists.
The relationship between data scale and linguistic competence is one of the
central empirical regularities of modern language modeling
\cite{kaplan2020scaling,hoffmann2022chinchilla}. This manuscript investigates a
complementary axis: structure.

\section{Background}
% P\=a\d{n}inian grammar as a generative system; Navya-Ny\=aya logic; structural
% priors and transfer in language models \cite{papadimitriou2020music}.

\section{Related Work}
% Compositional generalization benchmarks \cite{lake2018scan,kim2020cogs,
% keysers2020cfq}; minimal-pair syntactic evaluation \cite{warstadt2020blimp};
% the sample-efficient pretraining setting \cite{warstadt2023babylm}. Honest,
% verifiable citations only.

\section{Hypotheses}
% H1 grammar prior; H2 Ny\=aya scaffold + synergy; H3 epistemic constraint.

\section{Data Engine}
% Sa\d{m}s\=adhan\={\i} wrapper; diversity/coverage; sandhi-aware tokenizer;
% k-Shuffle Dyck control; corpus assembly and licensing.

The structural prior under test (H1) is supplied by a P\=a\d{n}inian
generator. Rather than depend on the non-redistributable Sa\d{m}s\=adhan\={\i}
toolset, we adopt the open-source Vidyut derivation engine, which produces, for
each generated form, the full s\=utra-by-s\=utra derivation as a gold structural
annotation. A preliminary diversity measurement over ti\d{n}anta forms derived
from twelve dh\=atus across the major ga\d{n}as yields a form type--token ratio
of $0.96$ with $132$ distinct s\=utras exercised and a mean derivation length of
$26.3$ steps, indicating that the generator does not collapse into degenerate
repetition and comfortably clears the pre-registered diversity threshold of
$0.80$; coverage scales further with the full Dh\=atup\=a\d{t}ha. Real-corpus
competence is grounded in the Digital Corpus of Sanskrit (Apache-2.0), whose
sentences exhibit bigram and trigram entropies of $0.98$ and $0.99$
respectively over a hundred-thousand-sentence sample. A sandhi-aware
tokenizer exposes morpheme boundaries recovered from the derivation, and a
k-Shuffle Dyck control is matched to the generator's surface statistics so that
any downstream advantage is attributable to grammatical structure rather than
to incidental distributional differences.

\section{Experimental Design}
% One model, three readings; arms A--G; size ladder; statistical protocol
% (bootstrap CIs, permutation tests, Holm--Bonferroni; pre-registered thresholds).

\subsection{H1 at proxy scale: arms, dose, and the readout}
The H1 test holds one decoder architecture fixed (a $\mu$P-parameterised
$100\,\mathrm{M}$-parameter model) and varies only the content of a bounded
structural pre-pretraining stream that precedes a common downstream phase on the
COGS compositional-generalisation benchmark~\cite{kim2020cogs}. Three arms isolate
the variable of interest: arm~A receives no structural pre-pretraining; arm~B
receives the P\=a\d{n}inian k\=araka-framed stream; and arm~C receives a
$k$-Shuffle Dyck stream whose surface statistics are matched to arm~B's, so that
any B-over-C difference is attributable to \emph{grammatical content} rather than
to generic hierarchical structure or incidental distributional differences. The
structural dose is equalised across B and C by a matched-epoch schedule over a
diversity-capped corpus. All three arms are run with three seeds.

A central methodological finding of this work is that the \emph{readout}, not the
prior, was the initial obstacle. Full logical-form exact-match generation floored
near $0.02$ for every arm, including arm~A, even at $150\,\mathrm{M}$ parameters
and a quadrupled step budget---an artefact of generation difficulty at proxy
scale, not an absence of compositional learning. We therefore moved to a
generation-free, teacher-forced \emph{role-discrimination} readout: the model is
scored as correct on an item when it assigns higher likelihood to the gold
logical form than to a minimally role-corrupted variant. To prevent
researcher degrees of freedom, the primary metric and the B-versus-C decision
thresholds were pre-registered blind for each readout in turn (ADR-0014 through
ADR-0016) and committed to version control before the corresponding comparison was
computed.

Final-state discrimination, however, \emph{saturated}: an arm-A-only calibration
sweep across four corruption tiers---an agent--theme filler swap and a
distractor-rebinding corruption, each on the full lexical-generalisation set and
on a non-canonical (passive and dative) subset where surface word order
misleads---placed the no-prior baseline between $0.91$ and $0.99$ on every tier.
With both endpoints saturated, the discriminating axis is \emph{sample
efficiency}: how quickly each arm climbs the role-competence curve. The
pre-registered primary metric (ADR-0016) is the normalised area under the
discrimination-versus-downstream-tokens curve (AUC), with tokens-to-threshold
($\theta = 0.80$) and an early-checkpoint accuracy gap as concordant secondaries
required to agree in direction; the effect bar is $\Delta_{\mathrm{AUC}} \ge 0.02$
with a bootstrap confidence interval excluding zero.

\section{Results}
% Filled from the experiment ledger at each phase closure. Figures live in
% paper/figures/. Tables report mean $\pm$ 95\% CI.

\subsection{H1: no advantage of the P\=a\d{n}inian prior over a matched control}
Table~\ref{tab:h1-se} reports the H1 sample-efficiency result. Both structural
pre-pretrainings (arms~B and~C) lift the discrimination-curve AUC by approximately
$0.015$ over the no-pretraining baseline (arm~A), but they are
statistically indistinguishable from each other: the primary contrast is
$\Delta_{\mathrm{AUC}}(\mathrm{B}-\mathrm{C}) = +0.003$ with a $95\%$ bootstrap
confidence interval of $[-0.006, +0.014]$, the early-checkpoint gap is $-0.011$
(non-concordant), and tokens-to-threshold shows no separation. The pre-registered
effect bar is not met, yielding a \textsc{no-or-weak-efficiency-gain} verdict. The
readout was demonstrably in a measurable regime: the baseline crossed
$\theta = 0.80$ within the $0.05$--$0.1$ budget region, so neither the
``curves-too-steep-to-resolve'' nor the ``not-learned-at-scale'' pre-registered
branch was triggered.

\begin{table}[t]
\centering
\caption{H1 sample-efficiency readout at $100\,\mathrm{M}$ parameters (COGS,
three seeds). AUC is the normalised area under the role-discrimination learning
curve; early accuracy is taken at $0.1$ of the token budget.}
\label{tab:h1-se}
\begin{tabular}{lcc}
\hline
Arm & Discrimination AUC & Early acc.\ ($0.1$ budget) \\
\hline
A (no pre-pretraining) & $0.888$ & $0.842$ \\
B (P\=a\d{n}inian) & $0.905$ & $0.835$ \\
C (Dyck control) & $0.902$ & $0.846$ \\
\hline
$\Delta$ (B$-$C) & $+0.003\ [-0.006, +0.014]$ & $-0.011$ \\
\hline
\end{tabular}
\end{table}

The outcome has the most informative possible shape, $\mathrm{A} <
(\mathrm{B} \approx \mathrm{C})$: the small benefit localises to \emph{generic
hierarchical structure}, which the matched Dyck control supplies equally, rather
than to P\=a\d{n}inian-specific annotation. Read together with the saturated
final-state readout, this closes the proxy-scale H1 question on both axes. We
report it as a clean controlled negative: at $100\,\mathrm{M}$ parameters, a
P\=a\d{n}inian structural prior confers no measurable advantage over a
surface-matched Dyck control for compositional argument-role acquisition, on
either final accuracy or sample efficiency. The most plausible explanation,
supported by the calibration sweep, is that the downstream data alone is
sufficient for argument-role composition at this scale---the baseline already
composes---so the prior has little room to add value on this particular
capability.

\section{Discussion}
The H1 negative is scoped, not universal: it concerns argument-role composition at
proxy scale, a capability the unprimed baseline already acquires from the
downstream data. Rather than weakening the broader thesis, the result sharpens it.
Local syntactic--semantic structure is precisely the segment of capability that
small models obtain cheaply; the place where explicit grammatical and relational
annotation could plausibly substitute for raw data volume is the segment the
scaling laws actually govern---factual and traced-reasoning capability. The live
claim therefore relocates from ``a grammatical prior aids composition'' to a
data-format hypothesis: that annotation synthesised by running the generator in
reverse over a structured knowledge base reaches matched factual coverage at fewer
training tokens. We treat that as a separate, falsifiable line of work rather than
as a rescue of the present null.

Two scope boundaries are stated explicitly to avoid over-reading the H1 negative.
First, H1 tests a P\=a\d{n}inian \emph{k\=araka} (L1) prior against the matched
control; it does \emph{not} test the typed Navya-Ny\=aya Paribh\=a\d{s}\=a (L2)
prior, which is the subject of a separate pre-registration (H1$'$). Second, the H1
readout is venue-saturated---the unprimed baseline already composes---so the result
is most precisely stated as ``this venue cannot adjudicate the prior on
argument-role composition,'' rather than as a universal absence of benefit. We
therefore make no claim here about Paribh\=a\d{s}\=a; preliminary H1$'$ numbers are
withdrawn as non-evidential pending a faithful, lossless L2 implementation and a
GPU-only, off-saturation measurement on official BLiMP/EWoK.

\section{Limitations}
% Grammar generates form, not facts; world-knowledge benchmarks reported only as
% reference points; single-accelerator fixed-budget framing; Z3 vy\=apti coverage
% limited to formally expressible rules.

\section{Conclusion}

\bibliographystyle{IEEEtran}
\bibliography{references}

\end{document}
